\section{考察}
\label{sec:discussion}

Sections~\ref{sec:selective-erosion}--\ref{sec:secondary} の正式な結果は、基礎となる技術そのものよりも制度的含意が直観的ではない単純なメカニズムを特定する。私的互換性が政治的に重要なのは、それが正式な排除の継続価値を変えるからである。企業は正式分割を変えることなく追加標準をサポートできるが、その私的な技術的応答は分割に付随する生産物市場rentを変え、それゆえ政府の正式な標準化レジームに対する誘因を変える。本稿の中心的貢献は、私的互換性が非互換な標準を単に低コスト化することではない。私的適応が、代替的な正式標準制度の相対的な政治価値を変え得ることである。

\subsection{技術適応から政治的誘因へ}
\label{subsec:discussion-mechanism}

メカニズムは、地域標準化連合の加盟国についての継続利得分解に最も明瞭に表れる。私的迂回がないとき、
\begin{equation}
    W_i^{SU,N}-W_i^{IS}
    =\mathscr E_i-\mathscr D_i,
    \label{eq:discussion-before}
\end{equation}
である。ここで $\mathscr E_i=m_i(K_M-K_I)+M_C(A-P)$ はcoalition加盟国市場で生じる加盟国政府の純排除余剰、$\mathscr D_i=m_o(P-C)$ は域外国市場における相互的不利益である。第1の対象は国民厚生項であり、純粋なproducer rentではない。それは国内消費者余剰の変化と、coalition加盟国の各仕向地における国内企業の営業利益差を組み合わせる。第2の対象は、同じ国内企業が、依然として標準が非互換な域外国市場に供給するときに負う利益上の不利益を捉える。

域外国企業のみの私的互換性は、これらの構成要素に非対称に作用する。ブロック標準を採用すると、域外国企業はcoalition加盟国市場で互換となる。加盟国の消費者余剰と国内企業利益のブロックはそこで完全互換のcounterpartへ移るため、coalition固有の構成要素 $\mathscr E_i$ は消える。しかしcoalition加盟企業は域外国標準を採用していない。そのため、域外国市場での企業の不利益は残り、$\mathscr D_i$ は存続する。迂回後の比較は
\begin{equation}
    W_i^{SU,O}-W_i^{IS}
    =-\mathscr D_i.
    \label{eq:discussion-after}
\end{equation}
である。中心的な変換はしたがって $\mathscr E_i-\mathscr D_i\rightarrow-\mathscr D_i$ である。

これが、侵食が完全ではなく選択的である理由である。私的互換性は地域標準化連合を国際標準化に変えることも、域外国を正式coalitionへ加入させることもない。分割は $\rho_{ij}^{SU}$ のままである。変わるのは、その分割が生み出す継続利得である。域外国企業の技術的応答は、coalition市場での排除を通じて加盟国を利していた地域的取り決めの部分を取り除く一方、ブロック外における加盟国の相互的不利益を残す。この意味で、私的な技術適応は、正式な制度境界そのものを直接変えることなく、その政治的価値を変える。

同じ会計は、Corollary~\ref{cor:technological-political} の「技術的代替、政治的補完」という解釈も説明する。私的互換性は、coalition加盟国市場におけるSUとISの実効互換性の差を小さくするという限定された意味で、多国間調和に対する技術的代替である。しかし政治的代替である必要はない。Section~\ref{sec:selective-erosion} より、
\[
    G_i^O-G_i^N=\mathscr E_i.
\]
である。$\mathscr E_i>0$ のとき、域外国企業のみの互換性は、多国間標準化へ移行する加盟国政府の相対利得を、私的迂回が取り除くrentと正確に同じだけ高める。差別的な地域的取り決めに固有の利得構成要素を侵食することにより、私的互換性はより広い調和を政治的に魅力的にし得る。これは特定されたmodel regionにおけるbinaryなgovernment-incentive comparisonであり、一般的なsupermodularityやcontinuous-action complementarityの主張ではない。

このfeedbackは、本稿のindustrial-organization channelでもある。正式標準はまず、生産物市場の互換関係と限界費用を形作る。次に企業は、追加標準をサポートすることによる営業利益増加が固定費用 $F$ を正当化するかを選ぶ。これらの選択が、各正式分割の下でCournot competition、利益、消費者余剰を変える。政府は正式分割それ自体ではなく、結果として得られる継続結果を評価する。したがって因果連鎖は、firm conduct $\rightarrow$ product-market outcomes $\rightarrow$ coalition-specific rents $\rightarrow$ government incentives over formal standards regimes である。私的互換性はすべての二国間frictionを外生的に一様に下げるものではない。採用はstrategicであり、標準固有であり、費用を伴い、現在の正式分割に依存する。

Theorem~\ref{thm:formal-coalition-stability} は、この企業レベルの応答が正式coalition stabilityへ到達する過程を示す。$F$ が高いとき、私的迂回は非活性である。地域加盟国は $\Omega_0$ 上で $W_M^{SU,N}>W^{IS}$ を成立させる排除優位を保持し、3つの対称な地域標準化連合すべてがstableである。$F$ が中間領域 $F_L<F<2T_A$ に低下すると、域外国企業だけが迂回を有利とする。加盟国利得は $W_M^{SU,O}<W^{IS}$ となり、域外国もISを厳格に選好する。grand coalitionはその結果すべての地域SUを厳格にblockし、維持されたstrict-blocking conceptの下で国際標準化が一意にstableとなる。したがって、特定された範囲では、より低い私的互換性費用は、coalition membershipを直接変えるのではなく継続利得を変えることによって、地域的な正式協力を不安定化し得る。

中間領域が経済的に有用なのは、結果がthreshold equalityに結び付いていないからである。Section~\ref{sec:coalition-stability} は、$F_L<2T_A$ を満たすproduct-market parameterのopen setを確立し、域外国企業のみの迂回を選ぶ非退化な互換性費用区間を残している。本稿はその区間外のmonotonicityを推論しない。互換性がさらに安価となり、より広範なmultistandarding patternが生じると、追加の採用結果が現れる。これらのlow-$F$ caseはheadline stability comparisonの外にあり、別に扱う。

\subsection{制度的解釈}
\label{subsec:discussion-institutions}

正式分割はstylized institutional arrangementとして読むのが適切である。各国別標準はfragmentation、地域標準化連合は一部の国の間のexclusive harmonization arrangement、国際標準化は3国すべてをカバーする共通標準を表す。これらのcategoryは、fragmented、exclusive、universalなcompatibility regimeの経済的区別を捉えるものであり、特定の標準化機関や協定を文字どおり記述することを意図しない。

私的互換性技術も同様の節度をもって解釈すべきである。$F$ を支払うことは、企業が追加的な正式標準をサポートできるようにする費用のかかる行動のreduced-form representationである。応用によっては、design change、multi-standard production、dual certification、interface adaptation、compliance engineering、その他のcompatibility layerに対応し得る。本モデルはこれらの技術を区別しない。メカニズムにとって重要なのは、投資が私的に選択され、離散的であり、特定の追加標準に結び付いているため、その収益がその標準のカバーする市場集合に依存することである。

したがって固定費用 $F$ には、単なる技術的bifurcation parameterを超える制度的解釈がある。$F$ が高いと、企業に正式標準の境界を迂回する誘因がないため、その境界は大きな経済的意味を持つ。中間領域では、域外国企業は2か国のブロック標準を有利にサポートできる一方、coalition加盟企業は域外国標準をなおサポートしない。この方向性のある応答こそが選択的侵食を生む。この結果は、互換性費用が下がればすべての費用水準で機械的に国際標準化が有利になる、という意味ではない。むしろ、正式調和の相対的な政治価値は、どの企業が迂回を有利とし、どのcoalition-specific rentがその投資によって取り除かれるかに依存する。

この区別は政策解釈にも節度を要求する。正式調和と私的互換性は、生産物市場で技術的代替であっても、政府の観点から政治的代替であるとは限らない。interoperabilityの費用や利用可能性の変化は正式な取り決め間のrentを再配分し、それによって地域的調和とより広い調和に対する政治的支持を変え得る。ただし、本モデルはそれらの取り決めに対するsocial-planner rankingを与えない。stabilityとefficiencyを混同してはならない。中間領域で $\rho^{IS}$ が一意にstableであることは、国際標準化がworld welfareを最大化することを意味せず、地域SUのinstabilityも地域調和が非効率であることを意味しない。本モデルの政府選好は、strict blockingを評価するために用いるnational-welfare comparisonである。

このメカニズムは、実証的検証それ自体を提供することなく、実証的な問いを示唆する。multi-standard production、interoperability investment、cross-standard certificationの費用低下が、地域的排他性への依存低下や、より広い調和に対する政府支持の変化と関連するかを問うことができる。このようなpatternの検証には、制度固有のmeasurementと別個のidentification strategyが必要である。現在のモデルはこれらの対象を結び付けるメカニズムを提供するが、その実証的な広がりや因果的な大きさを示すものではない。

\subsection{メカニズムの範囲とrobustness}
\label{subsec:discussion-scope}

対称的なMain Modelが有用なのは、stability mechanismをpartner selectionから切り離すからである。Theorem~\ref{thm:formal-coalition-stability} は市場規模の非対称性を必要としない。$m_1=m_2=m_3=1$ で、私的互換性によってregional SUの集合からISへ遷移するのに十分である。Section~\ref{sec:secondary} では、no-adoption continuation stateにおいて政府がどのregional partnerを好むかを明確化するためだけに $m_3=1-\delta$ を導入する。したがってProposition~\ref{prop:market-size-partner-selection} はcoalition stabilityを確立せず、市場規模の異質性が選択的侵食メカニズムを生むわけでもない。

モデル内部のロジックは4つの要素を示している。第1に、exclusive regional arrangementが一部市場で加盟国固有の継続優位を作る。第2に、私的な技術的応答がその優位を取り除くことができる。第3に、より広い取り決めの外に残ることに伴う何らかの不利益がその応答後も残る。第4に、政府が私的適応後に生じる継続利得を用いて正式制度を比較する。Cournot microfoundationではこれらの構成要素が $\mathscr E_i$ と $\mathscr D_i$ として明示される。この会計は、制度rentの一様ではなく選択的な侵食に基づくメカニズムを示唆するが、本稿は任意のoligopoly structureの下で同じ順位逆転が成立すると証明するものではない。

したがって、いくつかの特徴は本質的な主張の一部ではない。このメカニズムに市場規模の非対称性は不要であり、対称的なhigh-$F$ resultは一意のregional partnerを選ばない。ネットワーク効果はcanonical demand systemに入り、定量的な利得ブロックへ影響するが、選択的侵食をネットワーク効果それ自体が必要なものと解釈すべきではない。逆に、ネットワーク効果が無関係であると示したわけでもない。それは検証済みopen parameter regionが構築されるproduct-market environmentの一部であり続ける。

他の特徴はmodel-specificであり、定理の射程を限定する。3か国、1国1企業、数量競争、線形需要、binaryなcompatibility adoption、外生的なformal coalition menu、national welfareを最大化する政府を仮定する。stability resultはexclusive-membership partition gameにおけるstrict blockingを用いる。consent-based accession、farsighted stability、coalition-proofness、bargaining-based formationなどの代替概念は分析しない。これらの制約は、私的採用段階とすべての関連blocking comparisonを正確に解くことを可能にするため有用であるが、検証済みrobustness resultとみなすべきではない。

3か国環境は意図的に最小である。fragmentation、exclusive regional harmonization、universal harmonizationを区別するには十分豊かでありながら、私的採用とcoalition blockingをtractableに保つ。国数を増やすと、coalition size、複数の競合bloc、nestedまたはoverlapping compatibility structure、場合によってはpath-dependent accessionを扱う必要がある。それらは私的採用誘因と許容されるcoalition deviationの集合の双方を変え得る。したがって、より大きなcoalition networkへのメカニズム拡張には、3か国証明の機械的複製ではなく、追加のcoalition-formation structureが必要である。

Cournot competitionも同様の目的を持つ。それはcompatibility decisionを政府の継続利得へ直接結び付けるclosed-form quantityとprofit blockを与える。現在の定理は、differentiated Bertrand competition、heterogeneous firms、endogenous entryの下で同じstability reversalを確立していない。これらは、異なるstrategic competitionの下で選択的侵食の会計が維持されるかを検証する自然な方向であるが、frozen modelの外にある。

\subsection{限界と今後の研究}
\label{subsec:discussion-limitations}

制度的応用に特に関係する制約がさらに2つある。第1に、互換性はbinaryであり、その固定費用は外生的である。実際には、interoperabilityは部分的であったり、quality dependentであったり、累積投資によって構築されたりし得る。またその費用はR\&D、蓄積されたtechnical knowledge、certification infrastructure、regulatory designに依存し得る。互換性技術を内生化すると、予想される制度結果から企業の投資誘因へのfeedbackが生じる可能性がある。この可能性は、$F$ をprimitiveとし採用意思決定だけを内生化する現在の結果とは別である。

第2に、本モデルはtransfer、side payment、coalition bargainingを除外する。したがってstrict blockingは各deviator自身のnational-welfare gainに依存する。transferを認めると、coalition全体のgainと個々のblocking incentiveを切り離し得るため、これはnormalizationではなく実質的なextensionである。同様に、本モデルはaccession negotiation、endogenous standards design、repeated interaction、dynamic investment、incomplete informationを分析しない。これらはいずれもcoalition-specific rentの分配や持続性を変え得る。

これらの限界はformal theoremの射程を限定するが、モデル内部で確立されたメカニズムを損なうものではない。私的適応は正式な排除の継続価値を変える。それが地域coalitionを支える加盟国市場の優位を選択的に侵食しながら域外国市場の相互的不利益を残すとき、政府の順位を逆転させ、strict-blocking stable setを変え得る。次節では、このメカニズムと、私的互換性と正式な標準協力の関係に対する含意を要約して結論とする。
